% Part: lambda-calculus
% Chapter: church-rosser
% Section: definitions-and-properties

\documentclass[../../../include/open-logic-section]{subfiles}

\begin{document}

\olfileid{lam}{cr}{dap}

\olsection{व्याख्या आणि गुणधर्म}

या प्रकरणात चर्च–रॉसर गुणधर्माची कल्पना आणि त्या
गुणधर्माचे काही सामान्य पैलू मांडू.

\begin{defn}[चर्च–रॉसर गुणधर्म, CR] पदांवरील संबंध $\xredone$
  पुढील अट पूर्ण करत असेल, तर त्याला \emph{चर्च–रॉसर गुणधर्म}
  आहे असे म्हणतात: $M \xredone P$ आणि $M \xredone Q$
  असतील, तर असे~$N$ अस्तित्वात असते की
  $P \xredone N$ आणि $Q \xredone N$.
\end{defn}

लॅम्डा कलनाकडे संगणनाचे एक प्रतिमान म्हणून पाहता येते.
त्यात सामान्य रूपातील पदे ``मूल्ये'' आहेत आणि पदांवरील
न्यूनीकरणसंबंध हे ``संगणननियम'' आहेत. चर्च–रॉसर गुणधर्म
असे सांगतो की संगणन पुढे नेण्याचे एकाहून अधिक मार्ग
असले, तरी निकाल मिळत असल्यास त्याचे मूल्य एकमेव असते.

प्राथमिक बीजगणितातील उदाहरण घ्या. $4 \times (1+2) + 3$
मोजण्याचे अनेक मार्ग आहेत. आधी $1+2$ चे~$3$ करून
$4
\times 3+3$ मिळवता येते. किंवा वितरणनियम वापरून
आधी $4 \times (1+2)$ चे $4 \times 1+4
\times 2+3$ असे रूप करता येते. या दोन्ही रूपांचे
पुढील न्यूनीकरण $12+3$ मध्ये होते.

$\xredone$ म्हणजे $\beta$-न्यूनीकरण घेतले, तर
चर्च–रॉसर गुणधर्माचा एक परिणाम लगेच दिसतो:
पदाचे सामान्य रूप असेल, तर ते एकमेव असते.
समजा $M$ चे $P$ आणि $Q$ मध्ये न्यूनीकरण होते,
आणि ही दोन्ही सामान्य रूपे आहेत. चर्च–रॉसर गुणधर्मामुळे
असे $N$ असते की $P$ आणि $Q$ या दोन्हींचे त्यात
न्यूनीकरण होते. पण गृहीतकानुसार $P$ आणि $Q$ ही
सामान्य रूपे असल्याने $P$ आणि $Q$ यांचे $N$ मध्ये
न्यूनीकरण फक्त तुच्छ न्यूनीकरणच असू शकते.
म्हणजे $P$, $Q$ आणि $N$ ही समान आहेत.
त्यामुळे पदाच्या सामान्य रूपाचा उल्लेख आपण
\emph{एकमेव} म्हणून करू शकतो.

लॅम्डा कलनाकडे संगणनाचे प्रतिमान म्हणून पाहिल्यास,
एखाद्या पदाचे सामान्य रूप हा त्या पदापासून सुरू झालेल्या
संगणनाचा ``अंतिम निकाल'' मानता येतो. वरील निष्कर्षानुसार
असा अंतिम निकाल अस्तित्वात असेल, तर तो एकच असतो.
उदाहरणार्थ, $4 \times (1+2)+3$ मोजल्याचा निकाल
एकच, म्हणजे~$15$, असतो.

\begin{thm} \ollabel{thm:str}
  संबंध $\xredone$ ला चर्च–रॉसर गुणधर्म असेल आणि
  $\xred$ हा $\xredone$ समाविष्ट करणारा सर्वांत लहान
  संक्रमक संबंध असेल, तर $\xred$ लाही चर्च–रॉसर
  गुणधर्म असतो.
\end{thm}

\begin{proof}
  समजा
  \begin{align*}
    M & \xredone P_1 \xredone \dots \xredone P_m \text{ आणि}\\
    M & \xredone Q_1 \xredone \dots \xredone Q_n.
  \end{align*}
  उंची $m + 1$ आणि रुंदी $n + 1$ असलेली पदांची
  $N$ जाळी रचून प्रमेय सिद्ध करू. $i$ व्या ओळीतील
  आणि $j$ व्या स्तंभातील पद $N_{i,j}$ असे दर्शवू.

  $N_{i,j} \xredone N_{i+1,j}$ आणि
  $N_{i,j} \xredone N_{i,j+1}$ होईल अशी $N$
  रचू. तिची व्याख्या पुढीलप्रमाणे:
  \begin{align*}
    N_{0,0} &= M \\
    N_{i,0} &= P_i && \text{जर } 1 \le i \le m \\
    N_{0,j} &= Q_j && \text{जर } 1 \le j \le n \\
  \intertext{आणि इतर प्रसंगी:}
    N_{i,j} &= R
  \end{align*}
  येथे $R$ हे असे पद आहे की $N_{i-1,j} \xredone R$
  आणि $N_{i,j-1}
  \xredone R$. $\xredone$ च्या चर्च–रॉसर गुणधर्मामुळे
  असे पद नेहमी अस्तित्वात असते.

  आता $N_{m,0} \xredone \dots \xredone N_{m,n}$
  आणि $N_{0,n}
  \xredone \dots \xredone N_{m,n}$ मिळते.
  येथे $N_{m,0}$ हे $P_m$ आणि $N_{0,n}$ हे $Q_n$ आहे.
  $\xred$ च्या व्याख्येनुसार प्रमेय सिद्ध होते.
\end{proof}

\end{document}
